\begin{problem}
    El examen médico de tolerancia a la glucosa consiste en medir la respuesta del cuerpo a la glucosa. Se inicia midiendo la concentración de la glucosa basal y luego se ingiere un líquido con una determinada concentración de glucosa para monitorear su concentración en el cuerpo al transcurrir las horas.

    En el siguiente gráfico se representa la concentración de glucosa en el cuerpo de un paciente al realizarse este examen.

    \begin{center}
    (imagen)
    \end{center}

    ¿Cuál de las siguientes afirmaciones se puede deducir del gráfico?

    \begin{enumerate}[label=\Alph*.]
        \item El paciente ingiere la glucosa aproximadamente una hora después de iniciar el examen.
        \item La concentración de glucosa se estabiliza en su nivel basal luego de aproximadamente \SI{4}{h} de ingerir glucosa.
        \item La concentración de glucosa basal del paciente es de \SI{110}{mg/dL}, aproximadamente.
        \item El punto más alto de la concentración de glucosa sucede \SI{3}{h} después de ingerir la glucosa.
    \end{enumerate}
\end{problem}
